% Authoritative LaTeX; edit this file, not the historical Markdown.
\chapter{Position: RoPE From First Principles}\label{chapter-9---position-rope-from-first-principles}
A decoder produces sensible output for a short prompt, then becomes
unreliable when requests are batched or a cached prefix is reused. Every
tensor still has the expected shape. Every number is finite. What
changed? Perhaps the engine used a batch-row index as a token position.
Perhaps it rotated channels using the wrong pairing convention. Perhaps
a key that already contained position information was rotated a second
time. None of these errors needs to crash.

Position is therefore not a decorative label attached to a vector. It is
an input to a numerical transformation whose convention must survive
projection, layout conversion, scheduling, and eventually persistent
state. This chapter builds that transformation from two-dimensional
geometry, then turns the geometry into an explicit, tested ownership
contract.

We implement \textbf{standard rotary position embedding}, or RoPE, for
one token's head tensors. The new endpoint is positioned Q and K, with V
unchanged. We do not yet compute attention scores, causal masks,
softmax, or a KV cache. The absence of those later mechanisms is useful:
when a rotation is wrong, there is no attention implementation behind
which the error can hide.

\section{Give position a precise place in the
pipeline}\label{give-position-a-precise-place-in-the-pipeline}

\EngineFigure{ch09-journey}{Raw query and key owners pass through position-dependent rotations; the value owner bypasses RoPE.}{fig:ch09-journey}

Chapter 8 produces raw query, key, and value activations from a
normalized residual. Equal residuals under equal weights produce equal
raw projections. RoPE adds a new dependency: a logical token position.
In the decoder convention we are building, it transforms Q and K after
projection and before their future comparison. It leaves V alone. This
is an architecture choice, not a rule that every operation named
positional embedding must obey. Meta's Llama reference performs complex
rotations on Q and K at this boundary.
\href{https://raw.githubusercontent.com/meta-llama/llama-models/main/models/llama3/model.py}{Meta
model source}

For one token at logical position \(p\), the activation owners represent

\begin{equation}
\mathbf{Q}_p\in\mathbb{R}^{H_q\times D_h},\qquad
\mathbf{K}_p,\mathbf{V}_p\in\mathbb{R}^{H_{kv}\times D_h}.
\end{equation}

Subscript \(p\) identifies the token; it does not say the raw projection
has already been rotated. A tilde will identify a positioned activation.
Head counts can differ, as in grouped-query attention, but the Q/K head
width and rotary convention must agree for the comparisons we will
derive. Rotation acts independently within each head; it neither mixes
heads nor reads a different token's activation.

The scalar implementation accepts a generic \(\mathbf{X}\in
\mathbb{R}^{H\times D_h}\), used once for Q and once for K. Storage is
F32. Phase construction and pair arithmetic use F64, followed by an F32
output cast. The fresh-output API accepts validated immutable strided
views; the in-place API requires an exclusive canonical owner. Shape is
unchanged in both cases. These are physical contracts beyond the
real-number equation.

{\def\LTcaptype{none} % do not increment counter
\begin{longtable}[]{@{}
  >{\raggedright\arraybackslash}p{(\linewidth - 4\tabcolsep) * \real{0.3333}}
  >{\raggedright\arraybackslash}p{(\linewidth - 4\tabcolsep) * \real{0.3333}}
  >{\raggedright\arraybackslash}p{(\linewidth - 4\tabcolsep) * \real{0.3333}}@{}}
\toprule\noalign{}
\begin{minipage}[b]{\linewidth}\raggedright
Symbol
\end{minipage} & \begin{minipage}[b]{\linewidth}\raggedright
Meaning
\end{minipage} & \begin{minipage}[b]{\linewidth}\raggedright
Unit or constraint
\end{minipage} \\
\midrule\noalign{}
\endhead
\bottomrule\noalign{}
\endlastfoot
\(H\) & Heads in this call & Positive count \\
\(D_h\) & Coordinates per head & Positive count \\
\(R\) & Rotated prefix width & Positive, even, \(R\le D_h\) \\
\(i\) & Coordinate-plane index & \(0\le i<R/2\) \\
\(p,n\) & Logical token positions & Token intervals from an origin \\
\(\theta\) & Frequency base & Positive finite, dimensionless \\
\(\omega_i\) & Angular frequency of plane \(i\) & Radians per token
interval \\
\(\phi_i(p)\) & Rotation angle & Radians \\
\end{longtable}
}

\section{Construct a rotation instead of memorizing
one}\label{construct-a-rotation-instead-of-memorizing-one}
\index{RoPE!rotation}

\EngineFigure{ch09-plane}{A coordinate plane shows a counterclockwise rotation, with basis-vector images and a numerical pair calculation.}{fig:ch09-plane}

Start with a plane whose horizontal coordinate is \(a\) and vertical
coordinate is \(b\). A positive angle turns counterclockwise. A unit
horizontal vector rotated through \(\phi\) ends at
\((\cos\phi,\sin\phi)\). A unit vertical vector starts a quarter-turn
ahead, so its image is \((-\sin\phi,\cos\phi)\). Every vector \((a,b)\)
is a weighted sum of those two basis vectors. Rotating that sum gives
the matrix and its scalar meaning:

\begin{equation}
\begin{aligned}
\mathbf{R}(\phi)&=
\begin{bmatrix}\cos\phi&-\sin\phi\\\sin\phi&\cos\phi\end{bmatrix},\\
\begin{bmatrix}a'\\b'\end{bmatrix}
&=\mathbf{R}(\phi)\begin{bmatrix}a\\b\end{bmatrix}
=\begin{bmatrix}a\cos\phi-b\sin\phi\\a\sin\phi+b\cos\phi\end{bmatrix}.
\end{aligned}
\end{equation}

This is a \(2\times2\) linear map, not an addition of a learned position
vector. It changes direction without changing real-arithmetic length.
Expanding both squares makes the cancellation visible:

\begin{equation}
(a')^2+(b')^2
=(a^2+b^2)(\cos^2\phi+\sin^2\phi)=a^2+b^2.
\end{equation}

The mixed terms cancel because they have opposite signs. This
calculation also explains a common implementation bug. If the code
writes \(a'\) before reading the old \(a\) for the second equation, it
no longer implements the map above. Copy both original coordinates into
local variables before either store. An in-place algorithm needs no full
activation copy, but it does need that two-value dependency to remain
intact.

At position zero, the first plane's angle is zero, so \([1,2]\) remains
\([1,2]\). At position one, with frequency one radian per token, the
same pair becomes

\begin{equation}
\begin{aligned}
a'&=\cos 1-2\sin 1\approx-1.142640,\\
b'&=\sin 1+2\cos 1\approx 1.922076.
\end{aligned}
\end{equation}

The squared norm is still approximately five. The angle is \textbf{one
radian}, not one degree and not one complete revolution. A hand
calculation in degrees can look like a small harmless discrepancy while
testing a different operator. The executable fixture checks these values
against the scalar Rust result.

RoFormer introduces RoPE as a multiplicative positional transformation
whose query/key inner products expose relative position. We use that
mechanism as the starting point; the following elementary derivations
and engine contracts make its conventions explicit for this
implementation. \href{https://arxiv.org/html/2104.09864v5}{Su et al.,
RoFormer}

\section{Build a head from independent frequency
planes}\label{build-a-head-from-independent-frequency-planes}

\EngineFigure{ch09-frequencies}{Two unit-circle clocks receive the same position but rotate at frequencies one and one tenth.}{fig:ch09-frequencies}

A head is not one enormous two-dimensional rotation. Its rotary prefix
is partitioned into \(R/2\) planes. For the standard schedule selected
here,

\begin{equation}
\omega_i=\theta^{-2i/R},\qquad
\phi_i(p)=p\omega_i,\qquad 0\le i<R/2.
\end{equation}

The base controls how frequency varies across planes. For \(\theta>1\),
later planes rotate more slowly. The first plane has \(\omega_0=1\)
regardless of base. A mathematically positive base below one is accepted
by our API but reverses this ordering; accepting the parameter does not
mean a checkpoint was trained with it. Very small bases can also make
phases unrepresentable, which the execution path rejects before
mutation.

For \(R=4\) and \(\theta=100\), the two frequencies are \(1\) and
\(0.1\). Their continuous rotation periods are \(2\pi\) and \(20\pi\)
token intervals. These periods are not assertions that integer token
positions recur exactly, since neither period is an integer. Multiple
clocks give a richer phase pattern than a single clock; they do not
constitute a proof that arbitrary lengths remain distinguishable under
finite precision or useful to a trained model.

For adjacent pairing, the whole-head map is block diagonal:

\begin{equation}
\mathcal{R}_p=
\operatorname{diag}\!\left(
\mathbf{R}(p\omega_0),\ldots,
\mathbf{R}(p\omega_{R/2-1}),\mathbf{I}_{D_h-R}\right)
\in\mathbb{R}^{D_h\times D_h}.
\end{equation}

If \(R=D_h\), omit the identity tail. This matrix describes semantics;
the implementation never allocates it. It stores two coefficients per
plane and applies the two scalar equations directly. Materializing a
dense matrix would turn linear work into unnecessary quadratic storage
and arithmetic.

Our four-channel input \([1,2,3,4]\) at position one becomes
approximately

\begin{equation}
[-1.142640,\;1.922076,\;2.585679,\;4.279517].
\end{equation}

The first two coordinates use angle one; the last two use angle \(0.1\).
Rotating the whole array with one sine/cosine pair would pass a norm
test and still be the wrong frequency schedule.

\begin{quote}
\textbf{ENGINEERING FAILURE} Chapter 8's two-channel heads cannot test
base sensitivity. With \(R=2\), there is only plane zero, whose
frequency is always one. Keep that historical fixture for composition,
but add a four-channel fixture before claiming that changing the base
has been tested.
\end{quote}

\section{Why an absolute position produces a relative
score}\label{why-an-absolute-position-produces-a-relative-score}

\EngineFigure{ch09-relative}{Two rotation rays reveal a relative phase gap; a four-channel numerical trace checks the key-minus-query sign.}{fig:ch09-relative}

The key identity is not simply that rotations preserve length. It is
that rotations at different positions compose predictably. The transpose
undoes a rotation, and consecutive rotations add angles:

\begin{equation}
\mathbf{R}(\alpha)^{\mathsf T}=\mathbf{R}(-\alpha),\qquad
\mathbf{R}(\alpha)\mathbf{R}(\beta)=\mathbf{R}(\alpha+\beta).
\end{equation}

The first equality follows by transposing the two off-diagonal sine
entries. The second follows by expanding the matrix product and applying
the sine and cosine angle-addition identities. Because the whole-head
map applies these identities independently to disjoint planes, they also
hold for \(\mathcal{R}\) under one fixed frequency and pairing contract.

Let \(\mathbf{q},\mathbf{k}\in\mathbb{R}^{D_h}\) be fixed raw head
vectors, with query position \(p\) and key position \(n\). Their
positioned versions are
\(\widetilde{\mathbf{q}}_p=\mathcal{R}_p\mathbf{q}\) and
\(\widetilde{\mathbf{k}}_n=\mathcal{R}_n\mathbf{k}\). Now move the
query's rotation across the inner product:

\begin{equation}
\begin{aligned}
\widetilde{\mathbf{q}}_p^{\mathsf T}\widetilde{\mathbf{k}}_n
&=(\mathcal{R}_p\mathbf{q})^{\mathsf T}
  (\mathcal{R}_n\mathbf{k})\\
&=\mathbf{q}^{\mathsf T}\mathcal{R}_p^{\mathsf T}
  \mathcal{R}_n\mathbf{k}\\
&=\mathbf{q}^{\mathsf T}\mathcal{R}_{n-p}\mathbf{k}.
\end{aligned}
\end{equation}

The explicit positional dependency is \textbf{key position minus query
position}. The sign matters for arbitrary vectors. In our fixture,
\(\mathbf{q}=[1,2,3,4]^{\mathsf T}\),
\(\mathbf{k}=[2,-1,1,3]^{\mathsf T}\), \(p=2\), \(n=5\), \(R=4\), and
\(\theta=100\). Both routes yield approximately \(13.558047\); the Rust
paths differ by about \(5.25\times10^{-7}\) because they round different
intermediate F32 vectors. Using \(p-n\) on the right gives a different
score, not a rounding variant of the same identity.

This proof has several precise consequences. At the same position, the
raw dot product is preserved. Adding the same offset to both positions
leaves the fixed-vector score unchanged. Applying position \(p\) twice
yields position \(2p\), not position \(p\) again. Applying \(-p\) is an
inverse in real arithmetic, though F32 casts make a round trip
approximate.

It also has limits. Raw activations in a full decoder may already depend
on context and position through earlier layers; holding them fixed is a
premise of this identity. A causal visibility rule can change which keys
are compared. Changing a scaling configuration can invalidate the
shared-frequency premise. Common-shift invariance is not permission to
reset new query positions while leaving previously rotated keys in the
old coordinate system.

Nor does the identity make similarity decrease monotonically with
distance. Choose a one-plane
\(\mathbf{q}=\mathbf{k}=[1,0]^{\mathsf T}\). The score is \(\cos(n-p)\),
which rises again after falling. Claims about distributions of
frequencies, learned behavior, or long-distance envelopes are different
from a pointwise monotonicity guarantee. A researcher should test the
precise claim being made; an engineer should not build a pruning rule
from a stronger claim that the algebra does not support.

\section{Pairing is part of the checkpoint
interface}\label{pairing-is-part-of-the-checkpoint-interface}

\EngineFigure{ch09-pairing}{Adjacent and split-half storage orders connect the same named plane coordinates through different index pairs.}{fig:ch09-pairing}

The phrase ``rotate pairs'' leaves an essential question unanswered:
which coordinates are partners? For a prefix of width \(R\), this
chapter supports two explicit choices:

\begin{equation}
\begin{aligned}
\text{Adjacent:}\quad &(a_i,b_i)=(x_{2i},x_{2i+1}),\\
\text{Split half:}\quad &(a_i,b_i)=(x_i,x_{i+R/2}).
\end{aligned}
\end{equation}

With four channels, adjacent order is \([a_0,b_0,a_1,b_1]\) and
split-half order is \([a_0,a_1,b_0,b_1]\). Let permutation matrix
\(\mathbf{P}\) perform this reordering. The equivalent split-half
operator is

\begin{equation}
\mathcal{R}^{\mathrm{split}}_p
=\mathbf{P}\mathcal{R}^{\mathrm{adj}}_p\mathbf{P}^{\mathsf T},\qquad
\mathcal{R}^{\mathrm{split}}_p(\mathbf{P}\mathbf{x})
=\mathbf{P}(\mathcal{R}^{\mathrm{adj}}_p\mathbf{x}).
\end{equation}

Equivalence requires the operand to be permuted too. Switching a mode
flag while leaving projected coordinates unchanged is generally wrong.
If a conversion changes head coordinate order, the corresponding Q/K
weight rows must produce that order consistently. Conversion may also
interact with quantization groups and packed storage; it is not
necessarily an inexpensive metadata edit at model load time.

Both wrong and right conventions preserve norms. Both can produce finite
arrays with identical shapes. Therefore norm, finiteness, and length
checks are necessary but insufficient. Our coordinate oracle checks
complete output vectors for both conventions, and a permutation test
verifies equivalence only when data and operator are transformed
together. The unpermuted results must differ on the asymmetric fixture.

An integer constant in a backend cannot be interpreted from its name
alone. The pinned GGML CPU path supports adjacent NORMAL pairing and
split-half NEOX-style pairing. Hermon's paged preview helper uses
adjacent coordinates. These are source-verified paths, not
interchangeable defaults for every model. The exact implementation
record appears in the case study below.

\section{Partial rotation is a declared policy, not an odd-size
accident}\label{partial-rotation-is-a-declared-policy-not-an-odd-size-accident}

\EngineFigure{ch09-partial}{A five-channel head rotates its first four coordinates and preserves the last coordinate, including negative-zero bits.}{fig:ch09-partial}

Our configuration requires \(0<R\le D_h\) and even \(R\). It does
\textbf{not} require an even head width when only a prefix is rotated.
Thus \(D_h=5,R=4\) is valid; the last coordinate follows the identity
map. Conversely, \(D_h=4,R=3\) is invalid because one coordinate in the
requested prefix would lack a partner. We reject \(R=0\) rather than
interpreting it as a hidden disable switch.

Frequencies use denominator \(R\), the rotated prefix width. This is a
selected contract, not a universal promise about all partial-RoPE
architectures. A model adapter must inspect its actual definition.
Confusing full head width with rotary width changes every nonzero-index
plane even when all accesses remain within bounds.

Trailing finite values retain their exact F32 bits in both APIs. We
still validate the tail for finiteness. This intentionally makes the API
a checked finite-activation boundary, not merely a transform that
ignores invalid data outside its prefix. Position zero validates the
complete input and phases, then returns the unchanged finite payload;
the fresh-output route still owns a separate allocation. Signed zero is
preserved instead of being perturbed by an unnecessary multiply-add
sequence.

\section{Turn geometry into a failure-safe ownership
contract}\label{turn-geometry-into-a-failure-safe-ownership-contract}

\EngineFigure{ch09-transaction}{Validation and an immutable arithmetic preflight precede the first write; commit retains the activation allocation.}{fig:ch09-transaction}

The executable implementation lives in
\href{https://github.com/hermonai/inside-the-llm-engine/blob/astra-visual-rewrite/code/mini-engine/crates/engine0/src/rope.rs}{rope.rs}.
Its two entry points separate borrowing from mutation:

\begin{verbatim}
pub fn rope_reference(
    input: &TensorView<'_>,
    position: i64,
    config: RopeConfig,
) -> Result<OwnedTensor, RopeError>

pub fn rope_in_place(
    output: &mut OwnedTensor,
    position: i64,
    config: RopeConfig,
) -> Result<(), RopeError>
\end{verbatim}

These are API signatures, not complete function bodies. The linked
source contains all validation and execution.
\texttt{\seqsplit{RopeConfig::try\_new}} validates the dimensions and
base; private fields prevent callers from constructing an invalid
configuration behind the checks. The execution path then verifies a
nonempty rank-two tensor of the configured width, the position
conversion policy, every input coordinate, every frequency and angle,
and every output pair's representability as F32.

The last check matters. A finite rotation can concentrate a pair's
length into one coordinate. Two large finite F32 inputs can therefore
produce a coordinate outside finite F32 range even though the
real-number norm is preserved. F64 arithmetic postpones that overflow
but cannot make the final F32 storage represent it. Rejecting only
non-finite inputs is insufficient.

The in-place implementation uses three phases. First it validates input
and prepares one F64 sine/cosine pair per plane. Second it dry-runs
every pair without writing. Third it commits using the already validated
shape and phase plan. A failure in the last head's last plane cannot
leave earlier heads partially rotated. Every returned error preserves
the entire payload bit for bit, including NaN payloads in a rejected
input.

This is transactional behavior at a local numerical API boundary, not a
database transaction or recovery from process termination. Allocator
failure follows Rust's \texttt{Vec} behavior rather than returning a
\texttt{RopeError}. The phase scratch is \(O(R)\); the API is not
allocation-free. A process abort or hardware fault is outside the
returned-error guarantee. Naming those limits is part of the contract,
not a weakness to conceal.

After preflight, commit reads both old coordinates, evaluates the two
scalar equations, and stores both outputs. The exclusive owner prevents
a concurrent safe-Rust borrower from observing intermediate writes. The
operation retains the activation's payload pointer. By contrast,
\texttt{rope\_reference} prepares from the immutable view, materializes
a fresh contiguous owner, then commits there. It accepts padded, offset,
or broadcast input views because Chapter 5's logical indexing resolves
their physical addresses.

\begin{quote}
\textbf{PROVE IT} Test a late overflow in the second head, not only an
invalid first value. Snapshot every F32 bit before the call. After the
named error, compare the entire snapshot. A local pair test cannot prove
whole-owner failure atomicity.
\end{quote}

\section{Position precision is a separate numerical
problem}\label{position-precision-is-a-separate-numerical-problem}

An integer position must eventually participate in floating-point phase
construction. F32 represents every integer only through \(2^{24}\). In
the usual round-to-nearest conversion, \(2^{24}\) and \(2^{24}+1\)
become the same F32 value. If phase construction starts from those
identical values, it has already lost the distinction between the
positions. Better trigonometric evaluation cannot recover a distinction
discarded before the call.

Our API accepts signed 64-bit positions but restricts magnitude to
\(2^{53}\), the contiguous exact-integer range of F64. It rejects
out-of-policy positions before conversion, including both integer
extremes. Negative positions are useful for inverse and algebraic tests;
this is not a claim that a serving request should accept negative token
positions.

Exact integer conversion does not guarantee exact phase multiplication.
Frequency approximation, multiplication rounding, trigonometric argument
reduction, coefficient rounding, and final activation rounding remain. A
useful local sensitivity result comes directly from plane geometry. If a
phase error is \(\delta\) radians and \(\mathbf{x}\in\mathbb{R}^2\),
then

\begin{equation}
\left\|\mathbf{R}(\phi+\delta)\mathbf{x}
      -\mathbf{R}(\phi)\mathbf{x}\right\|_2
=2\|\mathbf{x}\|_2\left|\sin\frac{\delta}{2}\right|
\le\|\mathbf{x}\|_2|\delta|.
\end{equation}

The equality is the chord length between two points on a circle of
radius \(\|\mathbf{x}\|_2\). The bound uses \(|\sin z|\le|z|\). It
isolates phase sensitivity, not the combined error of every
floating-point operation. A phase error can move a vector substantially
while preserving its norm almost perfectly. Again, norm preservation
alone is not a sufficient quality gate.

Tests at positions \(1,7,2048,32768\) and \(-7\) exercise moderate-range
inverse and norm behavior with explicit F32 tolerances. Acceptance at
\(\pm2^{53}\) checks only that the conversion policy and finite-output
path hold. It does not certify accurate adjacent-position phases there,
let alone trained-model quality at that context length. Context
capacity, numerical distinguishability, and useful model behavior are
three different claims.

\section{Follow a token through the executable
composition}\label{follow-a-token-through-the-executable-composition}

Run the trace from the mini-engine workspace:

\begin{verbatim}
cargo run --quiet --example chapter09_rope_trace
\end{verbatim}

The example reuses Chapter 8's nonidentity projection fixture. It
selects embedding row one, normalizes with the learned gain, computes
three separate projections, transfers ownership of their head tensors,
and rotates Q and K at logical position seven. There are two heads of
width two, so only the first frequency exists. The resulting flattened
owners are approximately:

{\def\LTcaptype{none} % do not increment counter
\begin{longtable}[]{@{}ll@{}}
\toprule\noalign{}
Owner & Four stored coordinates after the new step \\
\midrule\noalign{}
\endhead
\bottomrule\noalign{}
\endlastfoot
Positioned Q & \([4.778205,1.742237,2.512963,2.674260]\) \\
Positioned K & \([0.176943,2.575917,0.310675,0.755081]\) \\
Original V & \([-0.365148,3.651481,0.182574,1.825741]\) \\
\end{longtable}
}

The separate four-channel fixture checks both frequencies and both
pairing modes. The Python oracle represents a pair as the complex number
\(a+\mathrm{i}b\) and multiplies by \(\exp(\mathrm{i}\phi)\). It does
not import Rust or reproduce the runtime's tensor indexing. It uses
independent real-arithmetic projection and normalization for the
preceding pipeline, so small F32 differences are expected rather than
hidden.

The
\href{https://github.com/hermonai/inside-the-llm-engine/blob/astra-visual-rewrite/scripts/check-rope-visual-parity.py}{parity
gate} compares all coordinates, both relative-dot routes, and the
complete composed Q/K/V against this oracle with \(10^{-5}+10^{-5}|y|\)
tolerance. This is a fixture acceptance criterion, not a universal error
theorem. The gate also deliberately alters and truncates outputs to
prove its comparator rejects bad candidates. The scene manifest binds
all ten figures to the same input-only fixture and requires each plate
to appear exactly once in this chapter.

\section{Logical coordinates must survive physical
rearrangement}\label{logical-coordinates-must-survive-physical-rearrangement}

\EngineFigure{ch09-positions}{Three requests have different logical positions, batch rows and physical block identifiers; only logical positions feed RoPE.}{fig:ch09-positions}

In a future batched engine, row zero may contain request A at position
127 while row one contains request B at position three. The row number
is an execution location; the token position is a sequence coordinate.
Sorting rows by length or priority must not silently change the latter.
Similarly, moving a stored key from one physical block to another
changes where its bytes live, not which rotation has already been
applied to them.

A multi-token prefill chunk beginning at logical position \(s\) would
rotate its token row \(t\) at \(s+t\), with checked integer addition. A
decode call would use the next logical position for that request. The
same per-token rotation law applies; scheduling determines which
positions enter the call. Our API accepts only one position and one
token's heads, so it does not implement this sequence wrapper or its
overflow checks by implication.

These observations define a future KV-state compatibility boundary. If
keys are stored after rotation, their interpretation depends on
positions, frequency schedule, pairing, rotary width, and model
parameters. Prefix reuse must preserve that interpretation. A matching
byte length or token string does not by itself prove compatible cached
activations. This is a derived design requirement, not a cache feature
shipped in this chapter.

Double rotation is particularly easy to disguise. An operation called
``prepare keys'' may rotate a fresh projection on one path and a
previously positioned key on another. Shapes and norms agree, but phase
now corresponds to the sum of both applications. Keep stage names
explicit: raw projection, positioned activation, and eventually stored
state. A future type-level state distinction can make accidental
repeated transformation harder to express.

\section{Long context changes geometry, not only a limit
field}\label{long-context-changes-geometry-not-only-a-limit-field}

\EngineFigure{ch09-scaling}{Analytical phase lines compare the original schedule, position interpolation and a changed base, with a first-plane counterexample to equivalence.}{fig:ch09-scaling}

Suppose the desired context is longer than the training context. Merely
allowing larger position integers exposes the model to a different range
of relative phases. It does not supply evidence that the model has
learned to use those phases. Context-extension methods address this
issue with modified geometry and empirical evaluation. Position
Interpolation studies rescaling positions into a shorter range; YaRN
develops a more elaborate extension method. Neither paper's reported
results automatically transfer to a custom inference implementation.
\href{https://arxiv.org/abs/2306.15595}{Position Interpolation},
\href{https://arxiv.org/abs/2309.00071}{YaRN}

The distinction between changing position and changing base is
algebraic. For an interpolation factor \(s>1\),

\begin{equation}
\phi_i^{\mathrm{interp}}(p)=\frac{p}{s}\theta^{-2i/R}.
\end{equation}

Every plane's phase slope is divided by \(s\). By contrast, replacing
base \(\theta\) with \(\theta'\) changes each frequency by

\begin{equation}
\frac{\omega_i'}{\omega_i}
=\left(\frac{\theta'}{\theta}\right)^{-2i/R}.
\end{equation}

That ratio depends on plane index. At \(i=0\) it is one: the first plane
does not change at all. Thus a base change cannot generally reproduce
uniform position interpolation. The plotted \(R=4\) example uses
\(\theta=100\), \(s=4\), and \(\theta'=10000\). The slower-plane slopes
are respectively \(0.1\), \(0.025\), and \(0.01\) radians per token. The
first-plane slopes are \(1\), \(0.25\), and \(1\).

Some real implementations use per-frequency factors, transition bands,
or amplitude adjustments. If coefficients include an amplitude
multiplier, unit-norm preservation is no longer the applicable invariant
without accounting for that multiplier. If a frequency schedule changes
with current sequence length, old and new keys cannot be assumed to
share a fixed rotation contract. Inspect the actual model configuration
and runtime behavior before reusing a proof derived for standard
fixed-frequency RoPE.

This chapter intentionally implements none of those scaling extensions.
Future support needs explicit configuration identity, checkpoint
compatibility, reference vectors, boundary tests around transition
regions, and task-level quality evaluation. A parameter accepted by a
loader is not the same thing as a model supported correctly.

\section{Price the implementation without inventing a speed
result}\label{price-the-implementation-without-inventing-a-speed-result}

With cached coefficients, one plane rotation uses four scalar multiplies
and two additions/subtractions. Counting each as one FLOP gives six
FLOPs per plane, hence \(3HR\) FLOPs per token tensor for the
transformation alone. Our preflight repeats that arithmetic before
commit, so it uses \(6HR\) FLOPs for pair evaluation at nonzero
positions. Frequency powers, trigonometric evaluation, checks, indexing,
allocation, and conversions are excluded from that count; those
exclusions can dominate the small scalar workload.

For F32 activation storage, a source-level access model for the in-place
nonzero path is

\begin{equation}
B_{\mathrm{activation}}
=4HD_h+4HR+8HR
=4HD_h+12HR\quad\text{bytes}.
\end{equation}

The first term scans all input coordinates for finiteness. The second
reads rotary coordinates in preflight. The last reads and writes those
coordinates in commit. This counts logical payload accesses, not
measured DRAM traffic. Cache reuse, compiler decisions, metadata checks,
and coefficient accesses are outside the model. The fresh-owner
reference adds a full activation copy and allocation. Its clearer
borrowing contract is useful as an oracle-facing path, not evidence of
performance parity with in-place execution.

The phase plan stores \(R/2\) pairs of F64 values, or \(8R\) payload
bytes, shared across heads of the token. A precomputed table for \(T\)
positions stores \(TR\) coefficients. At \(s_c\) bytes per coefficient,

\begin{equation}
B_{\mathrm{table}}=TRs_c\quad\text{bytes}.
\end{equation}

For \(T=32768,R=128,s_c=4\), that is 16 MiB, not counting metadata. A
table can be shared across heads, and potentially layers, only when
their complete frequency contracts agree. It is a phase table, not a KV
cache: it contains no token-dependent projected activations.
Lower-precision coefficients reduce table bytes but introduce different
phase and norm errors.

Direct computation pays for powers and trigonometric evaluation. A table
pays build time, memory footprint, indexing, and coefficient traffic.
Recurrence can update coefficients using angle addition but accumulates
drift and needs a reseeding policy. Fusion can avoid intermediate
activation traffic but complicates ownership, rounding, and the boundary
available for debugging. None of these candidates is implemented or
benchmarked here.

A valid future table-versus-compute experiment must compare identical
positions, \(H,D_h,R\), pairing, base, precision, and output tolerance.
Separate table construction from warm lookup; report cold-start cost
too. Declare allocation and preflight policies, contiguous versus
strided inputs, timing boundary, hardware, compiler, repetitions, and
raw results. Measure complete Q/K preparation as a second boundary
before claiming an end-to-end benefit. There is no reason to multiply a
kernel ratio into a serving throughput claim.

\section{Inside Hermon: follow the actual execution
route}\label{inside-hermon-follow-the-actual-execution-route}

\EngineFigure{ch09-source}{Default Batched execution, opt-in paged preview and the teaching implementation have distinct precision and configuration contracts.}{fig:ch09-source}

\begin{quote}
\textbf{INSIDE HERMON - CURRENT / PREVIEW / LIBRARY} Source inspected at
Hermon \texttt{2a3fd521} and its pinned llama.cpp \texttt{389ff61d} on
2026-09-12. Unset runtime mode selects Batched; paged remains opt-in
preview. A helper's existence is not evidence that the default route
calls it.
\end{quote}

The default route delegates model execution to llama.cpp. Its Llama
graph builds Q/K/V, then invokes \texttt{ggml\_rope\_ext} separately for
Q and K with the position input, rotary dimensions, mode, and
model/runtime scaling parameters. The pinned CPU kernel supports
multiple pairing modes, prepares phase terms for reuse across a token's
heads, and copies trailing non-rotary channels. Its extended path is
broader than the standard schedule implemented here.

Hermon's paged preview has a distinct \texttt{rope\_inplace} helper. It
iterates token rows and heads, pairs adjacent coordinates, computes
phases in F32, and optionally divides each frequency by a loaded factor.
Model construction checks the positive finite base, even head width, and
factor length/values. The helper also contains debug assertions; these
are not the same guarantee as a public API that returns typed errors in
a release build. The forward path rotates Q and K before subsequent KV
writes and attention.

The teaching implementation intentionally differs: F64 phases and pair
arithmetic, explicit partial-prefix policy, two pairing modes, and
preflight failure atomicity. It is not a claim of bitwise equality with
either production path. Source tests for scrambled physical block tables
and incremental decode illustrate why logical position must remain
separate from placement, but no real-model equivalence or hardware
performance run was performed for this milestone. The
\href{https://github.com/hermonai/inside-the-llm-engine/blob/astra-visual-rewrite/research/part-02/chapter-09-position-rope.md}{pinned
source record} identifies the exact paths and evidence classifications.

\section{Test the invariants that a plausible picture can
miss}\label{test-the-invariants-that-a-plausible-picture-can-miss}

The
\href{https://github.com/hermonai/inside-the-llm-engine/blob/astra-visual-rewrite/code/mini-engine/crates/engine0/tests/rope.rs}{focused
Rust tests} cover hand coordinates, both pairings, norm and inverse
behavior, rotation composition, the relative-score sign, common shifts,
exact tail bits, and fresh-versus-in-place ownership. Invalid
dimensions, bases, shapes, positions, non-finite inputs, unrepresentable
phases, and late F32 overflow exercise the failure contract.
Offset/strided and broadcast views separate mathematical coordinates
from contiguous storage assumptions.

Each test has a counterexample in mind. A norm test misses wrong pairing
and wrong frequency. An identity test at position zero misses nearly all
phase bugs. A first-head test misses bad head strides. A finite-input
test misses output overflow. A successful inverse does not prove useful
long-context quality. A same-buffer comparison cannot establish oracle
independence.

The implementation is safe Rust with no new dependencies. It remains
scalar and deliberately repeats arithmetic to make whole-owner failure
atomicity easy to prove. An optimized candidate may choose a different
policy, but it must name that policy rather than silently removing a
guarantee on which callers rely.

\section{Exercises: geometry, coordinates, and failure
boundaries}\label{exercises-geometry-coordinates-and-failure-boundaries}

The
\href{https://github.com/hermonai/inside-the-llm-engine/blob/astra-visual-rewrite/labs/lab-49-rope-position-workbench.md}{RoPE
workbench} specifies Labs 49-58, including inputs, expected artifacts,
deliberate breaks, and cleanup. Begin with the small numbers before
proposing a faster kernel.

\begin{enumerate}
\def\labelenumi{\arabic{enumi}.}
\tightlist
\item
  \textbf{CHECK:} Rotate \([1,2]\) at positions zero and one by hand.
  Explain the sign of each term and prove squared-norm preservation by
  expansion.
\item
  \textbf{CHECK:} Derive the relative-score identity with explicit
  transposes. Use asymmetric Q/K to make the reversed-sign version fail.
\item
  \textbf{BUILD:} Implement split-half pairing and its explicit
  permutation test. Show why a norm-only suite would accept the wrong
  coordinate order.
\item
  \textbf{BREAK:} Overwrite the first coordinate before calculating the
  second. Record which coordinate, norm, and inverse tests detect the
  mutation.
\item
  \textbf{CHECK:} Compare the old width-two fixture with the width-four
  fixture under a changed base. Identify the plane that cannot change.
\item
  \textbf{BREAK:} Place F32 maximum values in a late pair and verify
  that the entire owner remains unchanged after output-overflow
  rejection.
\item
  \textbf{BUILD:} Pass a padded or broadcast view to the fresh-output
  API and compare every coordinate with a contiguous input of equal
  logical values.
\item
  \textbf{EXTEND:} Design a phase-table key and a raw/positioned
  activation type boundary. State which mismatches must fail before
  cached state is reused.
\item
  \textbf{CHECK:} Derive phase-error sensitivity from a circle chord.
  Explain why nearly perfect norms can coexist with a large coordinate
  error.
\item
  \textbf{EXTEND:} Specify a reproducible table-versus-compute
  experiment. Keep quality evaluation separate from kernel accuracy and
  timing.
\end{enumerate}

\section{What this chapter establishes, and what comes
next}\label{what-this-chapter-establishes-and-what-comes-next}

Position now enters the executable engine through a checked numerical
API. Its meaning is visible in the basis vectors, frequency schedule,
relative inner-product identity, channel permutation, logical position,
and ownership transition. The implementation handles one token's heads;
it cannot silently substitute batching or cache placement for position
because neither exists inside the operator.

The next chapter consumes positioned Q and K to construct causal
attention. It must explain score scaling, visibility, stable softmax,
and value mixing without weakening the layout or precision contracts
established here. We have built the coordinates of comparison. We have
not yet built the comparison engine or the service that schedules it.

\subsection{Primary references and executable
evidence}\label{primary-references-and-executable-evidence}

\begin{itemize}
\tightlist
\item
  Su et al., \href{https://arxiv.org/html/2104.09864v5}{RoFormer}:
  rotary positional mechanism and its relative-position motivation.
\item
  \href{https://raw.githubusercontent.com/meta-llama/llama-models/main/models/llama3/model.py}{Meta
  Llama reference model}: complex Q/K rotation and architecture-specific
  frequency handling.
\item
  Chen et al., \href{https://arxiv.org/abs/2306.15595}{Position
  Interpolation}, and Peng et al.,
  \href{https://arxiv.org/abs/2309.00071}{YaRN}: context-extension
  research, not measured results of this teaching engine.
\item
  \href{https://github.com/hermonai/inside-the-llm-engine/blob/astra-visual-rewrite/research/part-02/chapter-09-position-rope.md}{Hermon
  source and review record},
  \href{https://github.com/hermonai/inside-the-llm-engine/blob/astra-visual-rewrite/code/reference/python/chapter09_rope_oracle.py}{complex
  oracle}, and
  \href{https://github.com/hermonai/inside-the-llm-engine/blob/astra-visual-rewrite/scripts/check-rope-visual-parity.py}{Rust/visual
  parity gate}.
\end{itemize}
\section{Worked synthesis: geometry must survive storage}
\subsection{Recover the relative sign}
\textbf{Problem.} For two-dimensional rotations, simplify
$(R(m)q)^\top(R(n)k)$. Which position difference remains?

\textbf{Solution.} Transpose reverses a rotation:
\begin{equation}
q^\top R(m)^\top R(n)k
=q^\top R(-m)R(n)k
=q^\top R(n-m)k.
\end{equation}
The remaining difference is key position minus query position in this
convention. Changing the sign while keeping the same rotation convention
changes the score. A norm-preservation test alone cannot detect this error.

\subsection{A norm-preserving wrong pairing}
\textbf{Problem.} For width four, compare adjacent pairs with split-half pairs.
Could the wrong choice pass a norm test?

\textbf{Solution.} Adjacent pairs are $(0,1),(2,3)$; split-half pairs are
$(0,2),(1,3)$. Each rotates disjoint coordinate planes and therefore preserves
the full norm in real arithmetic. Yet they generally produce different vectors.
Check component values against an independent oracle using the model's exact
pairing convention, not just an invariant both implementations share.

\subsection{Reject before mutating}
\textbf{Problem.} The first pair is valid, but rotating the last pair would
produce a nonfinite F32 output. What must a failure-atomic in-place API do?

\textbf{Solution.} Validate and dry-run all affected pairs before any write.
Return the typed error with every input bit unchanged. During the commit pass,
read both old coordinates before overwriting either one. This contract concerns
handled numerical errors; it does not promise recovery from allocator aborts.

